% !TEX root = ../aa_SoM_Chapters.tex

%% HARRIS: Chapter 9
%\xdef\Isectiontitle{Statistical Mechanics} 
%%
%\xdef\IaSubsectiontitle{A Simple Thermodynamic System} 
%\xdef\IbSubsectiontitle{Entropy and Temperature}
%\xdef\IcSubsectiontitle{Einstein's solid}
%\xdef\IdSubsectiontitle{Boltzmann \& Maxwell–Boltzmann distributions}
%\xdef\IeSubsectiontitle{Classical Averages}
%
%\xdef\IfSubsectiontitle{Quantum Distributions} 
%\xdef\IgSubsectiontitle{The Quantum Gas} 
%\xdef\IhSubsectiontitle{Photon Gas}
%\xdef\IiSubsectiontitle{The Laser}
%\xdef\IjSubsectiontitle{Specific Heats} 
%\xdef\IkSubsectiontitle{Debye Model} 

% ö = \%c3\%b6
% ä = \%C3\%A4

\xdef\sectiontitle{\IIIsectiontitle} %aiheuttama

%%%%%%%%%%%%%%%
\begin{frame}[label=3_4_a]%[label=3_4_TOC1]
\frametitle{\texorpdfstring{\culine[\sectiontitleulinecolor]{\sectiontitle}}{\sectiontitle} }

\vspace{1cm}

\begin{itemize}[leftmargin=0.5cm,itemsep=3mm]
    \item[3.1] \hyperlink{3_1_a}{\IIIaSubsectiontitle}
    \item[3.2] \hyperlink{3_2_a}{\IIIbSubsectiontitle}	
    \item[3.3] \hyperlink{3_3_a}{\IIIcSubsectiontitle}	
    \item[3.4] \hyperlink{3_4_a}{\CDAlert[red]{\IIIdSubsectiontitle}}
    \item[3.5] \hyperlink{3_5_a}{\IIIeSubsectiontitle}
    \item[3.6] \hyperlink{3_6_a}{\IIIfSubsectiontitle}
    \item[3.7] \hyperlink{3_7_a}{\IIIgSubsectiontitle}
\end{itemize}

%\endofslide 
\end{frame}
%%%%%%%%%%%%%%%

\xdef\sectiontitle{\IIIdSubsectiontitle} 

%%%%%%%%%%%%%%%
\begin{frame}
\frametitle{\texorpdfstring{\culine[\sectiontitleulinecolor]{\sectiontitle}}{\sectiontitle} }
\framesubtitle{Radioactive Decay processes}

\semismall
%On Radioactive Decay
\begin{itemize}
	\item Roughly \CDAlert[fuchsia]{3000 nuclide isotopes are known}\appear{, \CDAlert[red]{most} of them \CDAlert[red]{are radioactive}} 
	
	\item An \CDAlert[fuchsia]{unstable} radioactive nuclide \appear{has a tendency to spontaneously emit a small energetic particle} 
	\item According to estimates\appear{, in addition to the known ones,} \appear{2000 more are to be found}\appear{, all metastable or radioactive in some ways}
	
	\item There are \CDAlert[blue]{several kinds of 'decay processes'}, e.g.
\begin{itemize}
	\item $\alpha$ emission ($\mathrm{He}^{2+}$-ion)
	\item $\beta^{-}$ emission (electron)
	\item $\gamma$ emission (energetic $1 \mathrm{\;keV}–10 \mathrm{\;MeV}$ electromagnetic radiation)
\end{itemize}

\item[] And some \CDAlert[tuniblue]{less common processes}:
\begin{itemize}
\item Proton emission
\item Neutron emission
\item Positron emission
\item Double-beta emission (two electrons)
\item Some combinations of these
\end{itemize}

\end{itemize}

%\xdef\loc{south east}
%\begin{tikzpicture}[remember picture,overlay] % anchor=south
%    \node (A) [yshift=-0.25*\figYshift,xshift=\figXshift, anchor=\loc] at (current page.\loc) {\includegraphics[width=7cm]{Figures_booklet/SOM_12_Nuclear_physics_figures/SOM_Nuclear_physics_Figure_1.png}}; % 8cm max height
%\end{tikzpicture}

\endofslide 
%\endofsection
\end{frame}
%%%%%%%%%%%%%%%

%%%%%%%%%%%%%%%
\begin{frame}
\frametitle{\texorpdfstring{\culine[\sectiontitleulinecolor]{\sectiontitle}}{\sectiontitle} }
\framesubtitle{Radioactive Decay processes}

% 6.5 cm = midway, 10 cm = 3/4 
%\vspace{2.5mm}
\begin{minipage}{8.5cm}
\begin{itemize}
	\item The \Alert[red]{semi-empirical binding energy} formula can be used to study the stability of each isotope 
	
	\item Verified by experiments, we can plot the binding energies\appear{, schematically shown in the figure} \appear{(z-axis scaled by $100 \mathrm{\;MeV}$)} 
	
	\item One can also derive \CDAlert[blue]{theoretical limits} of how many protons or neutrons one can have in the nucleus 
	\implies{ These estimates for $N-Z$ boundaries \\
		generate the 'driplines' shown in lower figure }
\end{itemize}
\end{minipage}

\xdef\loc{south east}
\begin{tikzpicture}[remember picture,overlay] % anchor=south
    \node (A) [yshift=-0.25*\figYshift,xshift=\figXshift, anchor=\loc] at (current page.\loc) {\includegraphics[height=8.5cm]{Figures_booklet/SOM_Tipler_Nuclear_physics_figures/SOM_Tipler_Nuclear_physics_figure_20.png}}; % 8cm max height
\end{tikzpicture}

\endofslide 
%\endofsection
\end{frame}
%%%%%%%%%%%%%%%


%%%%%%%%%%%%%%%
\begin{frame}
\frametitle{\texorpdfstring{\culine[\sectiontitleulinecolor]{\sectiontitle}}{\sectiontitle} }
\framesubtitle{Radioactive Decay processes}

\seminormal
% 6.5 cm = midway, 10 cm = 3/4 
%\vspace{2.5mm}
\begin{minipage}{8.5cm}
\begin{itemize}
	\item Radioactive processes follow physical \CDAlert[magenta]{conservation laws}, conserved quantities being:
\begin{itemize}
	\item Relativistic mass–energy
	\item Electric charge
	\item Linear momentum
	\item Angular momentum
\end{itemize}

\item And also:
\begin{itemize}
	\item Nucleon number
	\item Lepton number
\end{itemize}

\item For a radioactive isotope, to be found in the nature \appear{(on Earth)}\appear{, the \Alert[red]{half-life} of the isotope must be shorter than the age of Earth} \appear{($4.5\times 10^{9}$ years)}\appear{, or the nuclide must be \CDAlert[blue]{continuously generated} by another radioactive substance}
\end{itemize}
\end{minipage}

\xdef\loc{south east}
\begin{tikzpicture}[remember picture,overlay] % anchor=south
    \node (A) [yshift=-0.25*\figYshift,xshift=\figXshift, anchor=\loc] at (current page.\loc) {\includegraphics[height=8.5cm]{Figures_booklet/SOM_Tipler_Nuclear_physics_figures/SOM_Tipler_Nuclear_physics_figure_20.png}}; % 8cm max height
\end{tikzpicture}

\endofslide 
%\endofsection
\end{frame}
%%%%%%%%%%%%%%%

%%%%%%%%%%%%%%%
\begin{frame}
\frametitle{\texorpdfstring{\culine[\sectiontitleulinecolor]{\sectiontitle}}{\sectiontitle} }
\framesubtitle{Alpha decay}
% 
\semismall
% 6.5 cm = midway, 10 cm = 3/4 
%\vspace{2.5mm}
\begin{minipage}{9.0cm}
	\begin{itemize}
	\item Alpha decay is a \Alert[red]{spontaneous emission of an energetic alpha particle}\appear{, i.e. He-4 nuclei.}
\appear{\xdef\loc{north east}%
\begin{tikzpicture}[remember picture,overlay]% anchor=south
    \node (A) [yshift=1*\figYshift,xshift=\figXshift, anchor=\loc] at (current page.\loc) {\includegraphics[width=4.7cm]{Figures_booklet/SOM_12_Nuclear_physics_figures/SOM_Nuclear_physics_Figure_18.png}}; % 8cm max height
\end{tikzpicture}%
} \appear{During the decay, an alpha particle is released out of the large nucleus} \appear{by tunneling through the potential energy barrier}
\appear{\xdef\loc{south east}%
\begin{tikzpicture}[remember picture,overlay] % anchor=south
    \node (A) [yshift=-0.25*\figYshift,xshift=\figXshift, anchor=\loc] at (current page.\loc) {\includegraphics[width=4.7cm]{Figures_booklet/SOM_Tipler_Nuclear_physics_figures/SOM_Tipler_Nuclear_physics_figure_1_cropped.png}};% 8cm max height
\end{tikzpicture}}

\item There is always reason for radioactive decay.
\appear{First of all, for the nucleus to be radioactive at all}\appear{, its \Alert[blue]{mass must be greater than the sum of the masses of the decay products}.} \appear{The energy released is given by}
\[
Q  \equal \appear{\textrm{((} m_{i}}  \minus \appear{m_{f}-}\appear{m_{\alpha} \textrm{))}} \appear{ c^{2}}
\]

\item The $\alpha$-active \CDAlert[fuchsia]{parent} nucleus is large \appear{(usually $A>200)$}\appear{, thus also the \CDAlert[fuchsia]{daughter} nucleus is rather heavy}\appear{, whereas the $\alpha$-particle is light: }\appear{\quad (\CDAlert[fuchsia]{conservation of momentum})}

\[
\appear{m \vec{v_d}}  \plus  \appear{M \vec{V_p}} \equal \appear{0} \quad \Rightarrow \quad \appear{v_d=}\frac{M}{m} \appear{V_p} \quad  \Rightarrow \quad  \frac{1}{2} \appear{m v_d^{2}}  \equal \frac{M}{m} \frac{1}{2} \appear{M V_p^{2}}
\]

\end{itemize}
\end{minipage}
%\vspace{2.5mm}

\endofslide 
%\endofsection
\end{frame}
%%%%%%%%%%%%%%%

%%%%%%%%%%%%%%%
\begin{frame}
\frametitle{\texorpdfstring{\culine[\sectiontitleulinecolor]{\sectiontitle}}{\sectiontitle} }
\framesubtitle{Alpha decay}

% 6.5 cm = midway, 10 cm = 3/4 
%\vspace{2.5mm}
\begin{minipage}{8.5cm}
\begin{itemize}
	\item Nucleons are bound by a strong nuclear force. \appear{However, sufficiently far from nucleus}\appear{, we are out of the range of the strong nuclear forces.} \appear{At that distance from core, there is a \CDAlert[blue]{positive Coulombic barrier}}
	
	\item Assumptions for explaining the $\alpha$ decay:
\end{itemize}
\end{minipage}
%\vspace{2.5mm}

\begin{itemize}
\item[] \begin{itemize}
		\item An $\alpha$-particle may exist as an entity within a heavy nucleus
	\item Such a particle is in constant motion \appear{and is held in the nucleus by a potential barrier}
	\item There is a small probability \appear{that the particle may tunnel through the barrier} \appear{with each collision with barrier}
\end{itemize}

\item Actually, it has also been suggested that\appear{, at least in the case of light nuclides}\appear{, the nucleons would be clustered into $\alpha$ particles...}
\end{itemize}

\xdef\loc{north east}
\begin{tikzpicture}[remember picture,overlay] % anchor=south
    \node (A) [yshift=1*\figYshift,xshift=\figXshift, anchor=\loc] at (current page.\loc) {\includegraphics[width=4.7cm]{Figures_booklet/SOM_Tipler_Nuclear_physics_figures/SOM_Tipler_Nuclear_physics_figure_25.png}}; % 8cm max height
\end{tikzpicture}

\endofslide 
%\endofsection
\end{frame}
%%%%%%%%%%%%%%%

%%%%%%%%%%%%%%%
\begin{frame}
\frametitle{\texorpdfstring{\culine[\sectiontitleulinecolor]{\sectiontitle}}{\sectiontitle} }
\framesubtitle{Alpha decay}

\begin{itemize}
	\item \CDAlert[red]{All heavy nuclei $(A > 210)$} \appear{\CDAlert[red]{are unstable} because the mass of the parent radioactive nucleus is greater} \appear{than the sum of the masses of decay products:} \appear{the daughter nucleus and an $\alpha$ particle}

\item In $\alpha$ emission for the nucleus \appear{$\Delta Z=\Delta N=-2$} \appear{thus $\Delta A=-4$}\appear{, and therefore there are four possible $\alpha$ decay chains} \vspace{3mm}
\[ 
\footnotesize
\begin{array}{r | c | c | c | c}
& \text {A} & \text {Longest lived isotope} & t_{1/2} & \text {Stable daughter} \\
\hline Thorium~series & 4 n  	   & n=58  & 1.4 \times 10^{10} \;a & Pb-208 \\
 Neptunium~series & 4 n+1 & n=58 & 2.1 \times 10^{6} \;a & Tl-205\\
 Uranium~series & 4 n+2 & n=59 & 4.5 \times 10^{9} \;a & Pb-206\\
 Actinium~series & 4 n+3 & n=58 & 7.0 \times 10^{8} \;a & Pb-207
\end{array}
\vspace{3mm}
\]

\item In each chain, the process proceeds via \CDAlert[blue]{subsequent $\alpha$ decays} towards lighter daughters. \appear{For the chain to follow the stability line}\appear{, \CDAlert[red]{also $\beta$ emissions must happen} every once in a while} \appear{(to get rid of extra neutrons)}

%\xdef\loc{south east}
%\begin{tikzpicture}[remember picture,overlay] % anchor=south
%    \node (A) [yshift=-0.25*\figYshift,xshift=\figXshift, anchor=\loc] at (current page.\loc) {$ \footnotesize
%\begin{array}{r | c | c | c | c}
%& \text {A} & \text {Longest lived isotope} & t_{1/2} & \text {Stable daughter} \\
%\hline Thorium~series & 4 n  	   & n=58  & 1.4 \times 10^{10} \;a & Pb-208 \\
% Neptunium~series & 4 n+1 & n=58 & 2.1 \times 10^{6} \;a & Tl-205\\
% Uranium~series & 4 n+2 & n=59 & 4.5 \times 10^{9} \;a & Pb-206\\
% Actinium~series & 4 n+3 & n=58 & 7.0 \times 10^{8} \;a & Pb-207
%\end{array}
%$}; % 8cm max height
%\end{tikzpicture}

\end{itemize}

%\xdef\loc{south east}
%\begin{tikzpicture}[remember picture,overlay] % anchor=south
%    \node (A) [yshift=-0.25*\figYshift,xshift=\figXshift, anchor=\loc] at (current page.\loc) {\includegraphics[width=7cm]{Figures_booklet/SOM_12_Nuclear_physics_figures/SOM_Nuclear_physics_Figure_1.png}}; % 8cm max height
%\end{tikzpicture}

\endofslide 
%\endofsection
\end{frame}
%%%%%%%%%%%%%%%

%%%%%%%%%%%%%%%
\begin{frame}
\frametitle{\texorpdfstring{\culine[\sectiontitleulinecolor]{\sectiontitle}}{\sectiontitle} }
\framesubtitle{Line of stability and decay chains}

% 6.5 cm = midway, 10 cm = 3/4 
%\vspace{2.5mm}
\begin{minipage}{9.0cm}
\begin{itemize}
	\item All series but one \CDAlert[fuchsia]{exist in nature}\appear{, except Neptunium-series (Np-237)}\appear{, because all the neptunium atoms have already decayed}

\item In the Figure we see the decay graph of Thorium series %\appear{and Uranium series}

\item The dashed line is the \CDAlert[blue]{line of stability}. \appear{The daughter nuclides of $\alpha$ decay are on the left side} \appear{, i.e. \CDAlert[fuchsia]{neutron rich side of the line}}\appear{, from where the process returns via sufficient numbers of $\beta$ decays}

\end{itemize}
\end{minipage}
\vspace{2.5mm}

\begin{itemize}
\item The \CDAlert[red]{radon gas} existing in room air \appear{(typical in Finland)} \appear{is Rn-222} \appear{which belongs to the Radium-series (U-238).} \appear{Its half-life is 4 days}\appear{, but more of it is continuously formed in the chain}
\end{itemize}


\xdef\loc{north east}
\begin{tikzpicture}[remember picture,overlay] % anchor=south
    \node (A) [yshift=1*\figYshift,xshift=\figXshift, anchor=\loc] at (current page.\loc) {\includegraphics[width=4.7cm]{Figures_booklet/SOM_Tipler_Nuclear_physics_figures/SOM_Tipler_Nuclear_physics_figure_26.png}}; % 8cm max height
\end{tikzpicture}


\endofslide 
%\endofsection
\end{frame}
%%%%%%%%%%%%%%%

%%%%%%%%%%%%%%%
\begin{frame}
\frametitle{\texorpdfstring{\culine[\sectiontitleulinecolor]{\sectiontitle}}{\sectiontitle} }
\framesubtitle{Decay chains}

\seminormal
% 6.5 cm = midway, 10 cm = 3/4 
%\vspace{2.5mm}
\begin{minipage}{9.0cm}
\begin{itemize}
	\item \CDAlert[blue]{Actinium-series}: \appear{The energy levels of Radium-223} \appear{after the $\alpha$ decay from Thorium-227}
	\appear{\xdef\loc{north east}%
\begin{tikzpicture}[remember picture,overlay] % anchor=south
    \node (A) [yshift=1*\figYshift,xshift=\figXshift, anchor=\loc] at (current page.\loc) {\includegraphics[width=4.7cm]{Figures_booklet/SOM_Tipler_Nuclear_physics_figures/SOM_Tipler_Nuclear_physics_figure_29.png}}; % 8cm max height
\end{tikzpicture}}
\item Also the corresponding energy spectrum of $\alpha$ radiation from Th-227 is shown \appear{(with the $\gamma$ energy given as a subscript)}
\appear{\xdef\loc{south west}%
\begin{tikzpicture}[remember picture,overlay] % anchor=south
    \node (A) [yshift=-0.25*\figYshift,xshift=-\figXshift, anchor=\loc] at (current page.\loc) {\includegraphics[width=8cm]{Figures_booklet/SOM_Tipler_Nuclear_physics_figures/SOM_Tipler_Nuclear_physics_figure_27.png}}; % 8cm max height
\end{tikzpicture}}
\item Note. Also here, as it very often happens\appear{, there will be a $\gamma$ quantum emitted in connection with the $\alpha$ emission}
\end{itemize}
\end{minipage}
%\vspace{2.5mm}



\endofslide 
%\endofsection
\end{frame}
%%%%%%%%%%%%%%%

%%%%%%%%%%%%%%%
\begin{frame}
\frametitle{\texorpdfstring{\culine[\sectiontitleulinecolor]{\sectiontitle}}{\sectiontitle} }
\framesubtitle{Beta decay}

% 
\begin{itemize}
	\item Electrostatic attraction cannot trap an electron into the nucleus\appear{, yet electrons are emitted by the nucleus}\appear{, being apparently spontaneously created in the decay process}

\item Note, \CDAlert[red]{in $\beta$ emission}\appear{, \CDAlert[red]{a neutron is converted into a proton}}\appear{, and thus the daughter nuclide will be a different element than the parent nuclide.} \appear{With a proton and an electron} 
\end{itemize}

% 6.5 cm = midway, 10 cm = 3/4 
%\vspace{2.5mm}
\begin{minipage}{8.2cm}
\begin{itemize}
	\item[] emitted\appear{, charge is conserved.} \appear{Energywise:}
$$
\begin{aligned}
\appear{Q} & \equal  \appear{\textrm{((} m_{n}-m_{p}-m_{e} \textrm{))}}  \appear{c^{2}} \\
& \equal \appear{(939.57-938.28-0.511) \mathrm{\;MeV}}  \equal \appear{0.78 \mathrm{\;MeV}}
\end{aligned}
$$
\appear{We assume this to be the kinetic energy of the product particles}\appear{, mostly of the electron}\appear{, since $m_{p} / m_{e}=1840$.} \appear{Empirically, the energy of the electron varies from nearly zero to $Q$}
\end{itemize}
\end{minipage}
%\vspace{2.5mm}


%% TABLE 1
%\xdef\loc{north east}
%\begin{tikzpicture}[remember picture,overlay] % anchor=south
%    \node (A) [yshift=0.9*\figYshift,xshift=\figXshift, anchor=\loc] at (current page.\loc) {\includegraphics[width=5.3cm]{Figures_booklet/SOM_12_Nuclear_physics_figures/SOM_Nuclear_Table_1.png}}; % 8cm max height
%\end{tikzpicture}


\xdef\loc{south east}
\begin{tikzpicture}[remember picture,overlay] % anchor=south
    \node (A) [yshift=-0.25*\figYshift,xshift=\figXshift, anchor=\loc] at (current page.\loc) {\includegraphics[width=5.3cm]{Figures_booklet/SOM_12_Nuclear_physics_figures/SOM_Nuclear_physics_Figure_19.png}}; % 8cm max height
\end{tikzpicture}

\endofslide 
%\endofsection
\end{frame}
%%%%%%%%%%%%%%%

%%%%%%%%%%%%%%%
\begin{frame}
\frametitle{\texorpdfstring{\culine[\sectiontitleulinecolor]{\sectiontitle}}{\sectiontitle} }
\framesubtitle{Beta decay}

%\seminormal
\begin{itemize}
	\item Historically, still three puzzles existed:
\begin{itemize}
	\item The fact that some of electrons have nearly zero energy \appear{(and \Alert[red]{not always $0.78 \mathrm{\;MeV}$})}\appear{, was impossible to explain by having only two product particles}
\end{itemize}
\end{itemize}

% 6.5 cm = midway, 10 cm = 3/4 
\vspace{2.8mm}
\begin{minipage}{8.5cm}
\begin{itemize}
\item[] \begin{itemize}[itemsep=2.8mm]
	\item \appear{Neutron}\appear{, proton} \appear{and electron}\appear{, are all spin 1/2-fermions}\appear{, so final total spin of $p^{+}$ and $e^{-}$ is an integer.} \appear{\Alert[blue]{Spin is not conserved?}}
	
\item Also, occasionally, \appear{the electron energy is nearly $Q$}\appear{, so if there is a third product particle}\appear{, it needs to have a \Alert[fuchsia]{very small mass}}
\end{itemize}

%\item Wolfgang Pauli suggested that there would\appear{, indeed, be a third product particle, a neutrino}\appear{, for which:}
%\begin{itemize}
%	\item spin would be $1 / 2$
%\item mass would be nearly zero \appear{(maybe $\textrm{((} \sim 2.2 \mathrm{ev} / \mathrm{c}^{2} \textrm{))} $}
%\item \appear{interaction with practically anything would be very weak}  \implies{ so-called \CDAlert[fuchsia]{weak interaction} }

%\end{itemize}
%\item In fact, the third particle $\bar{v}_{e}$ \appear{is referred as an antineutrino}
\end{itemize}
\end{minipage}
%\vspace{2.5mm}

\xdef\loc{south east}
\begin{tikzpicture}[remember picture,overlay] % anchor=south
    \node (A) [yshift=-0.25*\figYshift,xshift=\figXshift, anchor=\loc] at (current page.\loc) {\includegraphics[width=5.3cm]{Figures_booklet/SOM_12_Nuclear_physics_figures/SOM_Nuclear_physics_Figure_20.png}}; % 8cm max height
\end{tikzpicture}

\endofslide 
%\endofsection
\end{frame}
%%%%%%%%%%%%%%%

%%%%%%%%%%%%%%%
\begin{frame}
\frametitle{\texorpdfstring{\culine[\sectiontitleulinecolor]{\sectiontitle}}{\sectiontitle} }
\framesubtitle{Beta decay and neutrino}

%\seminormal
\begin{itemize}
	\item Wolfgang Pauli suggested that there would\appear{, indeed, be a third product particle, a \CDAlert[red]{neutrino}}\appear{, for which:}
\begin{itemize}
	\item spin would be $1 / 2$

\item mass would be nearly zero \appear{(maybe $\textrm{((}\sim 2.2 \mathrm{\;eV} / \mathrm{c}^{2} \textrm{))} $}

\item \appear{interaction with practically anything would be very weak}  \implies{ so-called \CDAlert[fuchsia]{weak interaction} }
\end{itemize}
\item In fact, the third particle $\bar{\nu}_{e}$ \appear{is referred \\
	to as an \CDAlert[blue]{antineutrino}}
\end{itemize}


\xdef\loc{south east}
\begin{tikzpicture}[remember picture,overlay] % anchor=south
    \node (A) [yshift=-0.25*\figYshift,xshift=\figXshift, anchor=\loc] at (current page.\loc) {\includegraphics[width=5.3cm]{Figures_booklet/SOM_12_Nuclear_physics_figures/SOM_Nuclear_physics_Figure_20.png}}; % 8cm max height
\end{tikzpicture}

\endofslide 
%\endofsection
\end{frame}
%%%%%%%%%%%%%%%

%%%%%%%%%%%%%%%
\begin{frame}
\frametitle{\texorpdfstring{\culine[\sectiontitleulinecolor]{\sectiontitle}}{\sectiontitle} }
\framesubtitle{Example of a $\beta^{-}$decay}

%
\begin{itemize}
	\item Let us consider the beta decay of boron-12:
\[
{ }_{5}^{12} B \quad \rightarrow \quad \appear{{ }_{6}^{12} C}  ~\plus~ \appear{{ }_{-1}^{~0} \beta^{-}}  ~\plus~ \appear{\bar{\nu}_{e}}
\]

\item Neutron is converted into proton\appear{, the new lighter nucleus is more tightly bound when $Z=N$} \appear{(and both are being even numbers).} \appear{Thus when boron decays, it will find lower energy}

\item Calculate the energies \appear{(we omit the mass of $\beta^{-}$}\appear{, since the atomic mass of ${ }_{6}^{12} C$} \appear{already contains one extra electron):}
$$
\begin{aligned}
\appear{Q} & \equal  \appear{\textrm{((} m_{\text {parent}}} \minus \appear{m_{\text {daughter }}}  \minus \appear{m_{\beta} \textrm{))}}  \appear{c^{2}}  \equal  \appear{\textrm{((} m_{B}-m_{C} \textrm{))}}  \appear{c^{2}} \\
&  \equal \appear{ \textrm{((} 12.014352 \;u-12 \;u \textrm{))}} \appear{\times 931.5 \;MeV / u} \equal \appear{13.4 \mathrm{\;MeV}}
\end{aligned}
$$
\end{itemize}

%\xdef\loc{south east}
%\begin{tikzpicture}[remember picture,overlay] % anchor=south
%    \node (A) [yshift=-0.25*\figYshift,xshift=\figXshift, anchor=\loc] at (current page.\loc) {\includegraphics[width=7cm]{Figures_booklet/SOM_12_Nuclear_physics_figures/SOM_Nuclear_physics_Figure_1.png}}; % 8cm max height
%\end{tikzpicture}

\endofslide 
%\endofsection
\end{frame}
%%%%%%%%%%%%%%%

%%%%%%%%%%%%%%%
\begin{frame}
\frametitle{\texorpdfstring{\culine[\sectiontitleulinecolor]{\sectiontitle}}{\sectiontitle} }
\framesubtitle{Beta decay of nitrogen-12}

\seminormal
\begin{itemize}
	\item Consider now $\beta^{+}$ decay of nitrogen-12\appear{ (${ }_{\,7}^{12} N$)}\appear{, during which a proton is converted into a neutron:}
\[
{ }_{7}^{12} N \quad \rightarrow \quad \appear{{ }_{6}^{12} C}  ~\plus~ \appear{{ }_{+1}^{\,0} \beta^{+}} \appear{~+~ \nu}
\]
\appear{where $\beta^{+}$ is a positive electron}\appear{, a positron}\appear{, taking care of the conservation of both charge and angular momentum}
\end{itemize}

% 6.5 cm = midway, 10 cm = 3/4 
\vspace{2.5mm}
\begin{minipage}{8.5cm}
\begin{itemize}
	\item[] \[ %what's this here
p^{+} \quad \rightarrow \quad \appear{n}  ~\plus~  \appear{e^{+}} \appear{~+~e}
\]
\item Calculate again $Q$ for the reaction:
$$
\begin{gathered}
Q = \textrm{((} m_{\text {parent}}-m_{\text {daughter}}-m_{\beta} \textrm{))}  c^{2}= \textrm{((} m_{N}-m_{C}-2 m_{e} \textrm{))}  c^{2} \\
= \textrm{((} 12.018613 \;u-12 \;u-2 \cdot 5.486 \times 10^{-5} \;u \textrm{))}  c^{2} =16.3 \mathrm{\;MeV}
\end{gathered}
$$

\item Note: A \Alert[blue]{free proton cannot decay into a neutron and a positron}, \appear{but a proton bound into a nucleus is able to do it}
\end{itemize}

\end{minipage}
%\vspace{2.5mm}

\xdef\loc{south east}
\begin{tikzpicture}[remember picture,overlay] % anchor=south
    \node (A) [yshift=-0.25*\figYshift,xshift=\figXshift, anchor=\loc] at (current page.\loc) {\includegraphics[width=4.6cm]{Figures_booklet/SOM_12_Nuclear_physics_figures/SOM_Nuclear_physics_Figure_21.png}}; % 8cm max height
\end{tikzpicture}

\endofslide 
%\endofsection
\end{frame}
%%%%%%%%%%%%%%%

%%%%%%%%%%%%%%%
\begin{frame}
\frametitle{\texorpdfstring{\culine[\sectiontitleulinecolor]{\sectiontitle}}{\sectiontitle} }
\framesubtitle{Electron capture}

%Electron capture
\begin{itemize}
	\item In the \CDAlert[red]{electron capture reaction}\appear{, proton in the nucleus captures an electron from the atomic electron shells} \appear{and changes into a neutron} 
	
	\item \CDAlert[red]{$Z$ decreases by one and $N$ increases by one}, just like in $\beta^{+}$ decay. \appear{In fact, many $\beta^{+}$-active nuclei}\appear{ can also capture electrons}

\item Always, when the mass of a nucleus with a given atomic number $Z$ is larger than the mass of a nucleus with atomic number $Z-1$\appear{, the electron capture will be possible}
\end{itemize}

% 6.5 cm = midway, 10 cm = 3/4 
\vspace{2.5mm}
\begin{minipage}{8.5cm}
\begin{itemize}
	\item If the mass difference is at least $2 m_{e} c^{2}$\appear{, then also a $\beta^{+}$ decay is possible.} \appear{In this case these \CDAlert[blue]{two decay processes compete}}

\item The probability of electron capture is usually negligible \appear{unless the atomic electron is in immediate vicinity of the nucleus}
\end{itemize}
\end{minipage}
%\vspace{2.5mm}


\xdef\loc{south east}
\begin{tikzpicture}[remember picture,overlay] % anchor=south
    \node (A) [yshift=-0.25*\figYshift,xshift=\figXshift, anchor=\loc] at (current page.\loc) {\includegraphics[width=4.8cm]{Figures_booklet/SOM_12_Nuclear_physics_figures/SOM_Nuclear_physics_Figure_22.png}}; % 8cm max height
\end{tikzpicture}

\endofslide 
%\endofsection
\end{frame}
%%%%%%%%%%%%%%%

%%%%%%%%%%%%%%%
\begin{frame}
\frametitle{\texorpdfstring{\culine[\sectiontitleulinecolor]{\sectiontitle}}{\sectiontitle} }
\framesubtitle{Example of an electron capture}

% 
\begin{itemize}
	\item Consider now electron capture reaction of carbon-11\appear{, ${ }_{6}^{11} C$}
\[
{ }_{6}^{11} C \appear{~+~} \appear{ _{-1}^{\;0} \beta^{-}} \quad \rightarrow \quad \appear{{ }_{5}^{11} B}  ~\plus~ \appear{\nu} \qquad \appear{\textrm{((}} \appear{p^{+}}  ~\plus~ \appear{e^{-}} \quad \rightarrow \quad  \appear{n~+~\nu} \appear{\textrm{))}}
\]
\appear{The energy released} \appear{will depend on the difference of the masses of parent and daughter:} 
\[
Q  \equal  \appear{\textrm{((} m_{\text {parent}}}  \minus \appear{m_{\text {daughter}} \textrm{))} } \appear{c^{2}}
\]

\item Another typical example of electron capture is 
\[
 { }_{24}^{51} C r  ~\plus~ \appear{{ }_{23}^{51} V}  ~\plus~ \appear{\nu_{e}} \appear{\,,~ \text{with}} \qquad  \appear{Q=}\appear{0.751 \mathrm{\;MeV}}
 \]
\end{itemize}

%\xdef\loc{south east}
%\begin{tikzpicture}[remember picture,overlay] % anchor=south
%    \node (A) [yshift=-0.25*\figYshift,xshift=\figXshift, anchor=\loc] at (current page.\loc) {\includegraphics[width=7cm]{Figures_booklet/SOM_12_Nuclear_physics_figures/SOM_Nuclear_physics_Figure_1.png}}; % 8cm max height
%\end{tikzpicture}

\endofslide 
%\endofsection
\end{frame}
%%%%%%%%%%%%%%%

%%%%%%%%%%%%%%%
\begin{frame}
\frametitle{\texorpdfstring{\culine[\sectiontitleulinecolor]{\sectiontitle}}{\sectiontitle} }
\framesubtitle{Decay chains}

\semismall
% 6.5 cm = midway, 10 cm = 3/4 
%\vspace{2.5mm}
\begin{minipage}{9.0cm}
\begin{itemize}[itemsep=1.5mm]
	\item Revisit the presentation of the thorium and uranium series of alpha decay as a $N Z$-graph

\item Note that to keep the isotopes in the vicinity of the \CDAlert[fuchsia]{line of stability}\appear{, there needs to be corrective beta decays}\appear{, especially if there's an abundance of neutrons}

\item Eventually, the \CDAlert[red]{decay pathways end} at around isotopes of \Alert[red]{lead}: \appear{Pb-208}\appear{, Pb-206} \appear{and Pb-207 for the thorium}\appear{, uranium} \appear{and actininium series, respectively.} \appear{It is not that those isotopes of lead would be universally very stable}\appear{, it is just that they are the stable isotopes \CDAlert[red]{'first in line'}}\appear{, and there are no further decay pathways form those isotopes.} \appear{(unlike certain distinct smaller isotopes e.g. K-40, C-14 etc.)}

\item Note also:
\begin{itemize}
	\item $\beta^{-}$ decay is directed diagonally down \& right 
	\item $\beta^{+}$ decay is directed diagonally up \& left
\end{itemize}

\end{itemize}
\end{minipage}
%\vspace{2.5mm}

\xdef\loc{east}
\begin{tikzpicture}[remember picture,overlay] % anchor=south
    \node (A) [yshift=0.*\figYshift,xshift=\figXshift, anchor=\loc] at (current page.\loc) {\includegraphics[height=8cm]{Figures_booklet/SOM_12_Nuclear_physics_figures/SOM_Nuclear_physics_Figure_23.png}}; % 8cm max height
\end{tikzpicture}

\endofslide 
%\endofsection
\end{frame}
%%%%%%%%%%%%%%%

%%%%%%%%%%%%%%%
\begin{frame}
\frametitle{\texorpdfstring{\culine[\sectiontitleulinecolor]{\sectiontitle}}{\sectiontitle} }
\framesubtitle{Valley of stability}

\semismall
% 6.5 cm = midway, 10 cm = 3/4 
%\vspace{2.5mm}
%\begin{minipage}{8.0cm}
\begin{itemize}[itemsep=1.6mm]
	\item One can also investigate the \CDAlert[blue]{stability valley} of nuclei\appear{, using $\beta$ decay processes}

\item With the $\beta$ decays \appear{we move in perpendicular direction to the energy valley}

\item The energy profiles of constant atomic mass show a parabolic type of energy valley

\xdef\loc{south east}
\begin{tikzpicture}[remember picture,overlay] % anchor=south
    \node (A) [yshift=-0.15*\figYshift,xshift=\figXshift, anchor=\loc] at (current page.\loc) {\includegraphics[width=11cm]{Figures_booklet/SOM_Tipler_Nuclear_physics_figures/SOM_Tipler_Nuclear_physics_figure_33.png}}; % 8cm max height
\end{tikzpicture}

\item The energy values here were \Alert[red]{computed by the semi-empirical binding energy formula}

\item For the odd $A$ series\appear{, the pairing term is zero \textrm{((}$C_{5}=0$\textrm{))}}\appear{, for the even $A$}\appear{, the $\beta$ decay}
\end{itemize}
% 6.5 cm = midway, 10 cm = 3/4 
%\vspace{1.4mm}
\begin{minipage}{2.8cm}
\begin{itemize}
	\item[] series jumps between two parabola\appear{, with $\pm C_{5}$.} \appear{(either $Z$-even/$N$-even or $Z$-odd/$N$-odd)}
\end{itemize}
\end{minipage}
%\vspace{2.5mm}

%\end{minipage}
%\vspace{2.5mm}



\endofslide 
%\endofsection
\end{frame}
%%%%%%%%%%%%%%%

%%%%%%%%%%%%%%%
\begin{frame}
\frametitle{\texorpdfstring{\culine[\sectiontitleulinecolor]{\sectiontitle}}{\sectiontitle} }
\framesubtitle{Gamma decay}

%Gamma decay
\begin{itemize}
	\item In $\gamma$ decay\appear{, a \CDAlert[fuchsia]{nucleus of an excited state decays} to lower energy level and emits a photon}\appear{, maintaining the original isotope.} \appear{The process is analogous to emitting a photon from atomic energy states}

\item \CDAlert[red]{ATOMIC}: For atomic processes the energies are in the range $\mathrm{\;eV}-\mathrm{\;MeV}$ \appear{and emitted wavelengths are $\mu \mathrm{m}-\mathrm{\;nm}$}

\item \CDAlert[blue]{NUCLEAR}: For nuclear processes the energies exceed MeV\appear{, and wavelengths are $\appear{\lambda} \equal \frac{h c}{E} \approx \appear{10^{-3} \mathrm{\;nm}}$
}
\item Note, that for a given quantum of electromagnetic radiation in the keV–MeV energy range\appear{, there is fundamentally really no way to distinguish}\appear{, whether it is an X-ray} \appear{or a $\gamma$-ray}\appear{, if\\ no further information of the system is given!}
\end{itemize}

\xdef\loc{south east}
\begin{tikzpicture}[remember picture,overlay] % anchor=south
    \node (A) [yshift=-0.25*\figYshift,xshift=\figXshift, anchor=\loc] at (current page.\loc) {\includegraphics[width=4.3cm]{Figures_booklet/SOM_12_Nuclear_physics_figures/SOM_Nuclear_physics_Figure_24.png}}; % 8cm max height
\end{tikzpicture}

\endofslide 
%\endofsection
\end{frame}
%%%%%%%%%%%%%%%

%%%%%%%%%%%%%%%
\begin{frame}
\frametitle{\texorpdfstring{\culine[\sectiontitleulinecolor]{\sectiontitle}}{\sectiontitle} }
\framesubtitle{Combined Alpha/beta and gamma decay}

\begin{itemize}
	\item After $\alpha$-, and $\beta$-decay\appear{, the \Alert[red]{nucleus is usually left into an excited state}}\appear{, and a $\gamma$ emission will follow} 
	
	\item The lifetimes of those excitation states are usually very short\appear{, unless selection rules restrict de-excitation of the excited states} 
	
	\item In practice, different $\gamma$ energies are often \Alert[red]{connected with specific decay processes}/shift in the nuclear energy levels

\item In general, $\gamma$ energies are characteristic for given isotopes\appear{, and they can be used as a \CDAlert[red]{fingerprint}}\appear{, i.e. to give information to identify the isotopes}\appear{, and also give further information e.g. of the energies at the shell structure of the atomic nucleus} \appear{(recall earlier discussion of Th-227 emission spectra)}
\end{itemize}

%\xdef\loc{south east}
%\begin{tikzpicture}[remember picture,overlay] % anchor=south
%    \node (A) [yshift=-0.25*\figYshift,xshift=\figXshift, anchor=\loc] at (current page.\loc) {\includegraphics[width=7cm]{Figures_booklet/SOM_12_Nuclear_physics_figures/SOM_Nuclear_physics_Figure_1.png}}; % 8cm max height
%\end{tikzpicture}

\endofslide 
%\endofsection
\end{frame}
%%%%%%%%%%%%%%%


%%%%%%%%%%%%%%%
\begin{frame}
\frametitle{\texorpdfstring{\culine[\sectiontitleulinecolor]{\sectiontitle}}{\sectiontitle} }
\framesubtitle{Internal conversion}

%Internal conversion
\begin{itemize}
	\item An alternative to $\gamma$-ray emission due to de-excitation of the excited nuclear state is \CDAlert[blue]{internal conversion}

\item During this conversion, the energy from the excited nucleus is \Alert[blue]{transferred into an} \CDAlert[blue]{orbital electron}\appear{, which will be emitted from the atom}\appear{ at high velocities}
 
 \item This concerns mostly \CDAlert[blue]{$s$-electrons of $K$- and $L$-shells}\appear{, which have a significant probability to reside close to the atomic core} 
 
 \item The energy will be given by:
\[
hf  \equal \appear{ \textrm{((} E_{\text {high}}}\appear{-E_{\text {low}} \textrm{))}} \appear{-B_{\text {electron}}}
\]

\item This kind of internal conversion is a one-step many-particle process\appear{, although in principle it can be considered to consist of two stages}\appear{, emission and absorption of $\gamma$}
\end{itemize}

%\xdef\loc{south east}
%\begin{tikzpicture}[remember picture,overlay] % anchor=south
%    \node (A) [yshift=-0.25*\figYshift,xshift=\figXshift, anchor=\loc] at (current page.\loc) {\includegraphics[width=7cm]{Figures_booklet/SOM_12_Nuclear_physics_figures/SOM_Nuclear_physics_Figure_1.png}}; % 8cm max height
%\end{tikzpicture}

\endofslide 
%\endofsection
\end{frame}
%%%%%%%%%%%%%%%

%%%%%%%%%%%%%%%
\begin{frame}
\frametitle{\texorpdfstring{\culine[\sectiontitleulinecolor]{\sectiontitle}}{\sectiontitle} }
\framesubtitle{Spontaneous fission}
\seminormal

\begin{minipage}{8.6cm}
\begin{itemize}
	\item Some heavy radioactive nuclei have a tendency to decay differently, \appear{i.e. to \Alert[red]{split into two larger fragments} by \CDAlert[red]{spontaneous fission}}

\item Driving force is \Alert[fuchsia]{minimization of total energy}

\item Two separate nuclei will have \Alert[red]{larger surface area}\appear{ than the initial nucleus} \appear{(and thus a higher energy)} 
\end{itemize}
\end{minipage}
\vspace{2.5mm}

\begin{itemize}
\item \appear{But the two separate nuclei can also \Alert[blue]{arrange the protons} of the nuclei further away from each other}
\implies{ \Alert[blue]{Coulombic repulsion decreases}\appear{, increasing binding forces\\} \appear{and compensating the increased surface energy} }

\item Another important detail in the spontaneous fission is\appear{, that the fragmentation will generate also \CDAlert[red]{free neutrons}}\appear{, which can induce further reactions in the system by colliding into the abundant parent nuclei}
\end{itemize}

\xdef\loc{north east}
\begin{tikzpicture}[remember picture,overlay] % anchor=south
    \node (A) [yshift=0.8*\figYshift,xshift=\figXshift, anchor=\loc] at (current page.\loc) {\includegraphics[width=5cm]{Figures_booklet/SOM_12_Nuclear_physics_figures/SOM_Nuclear_physics_Figure_25.png}}; % 8cm max height
\end{tikzpicture}

%\endofslide 
\endofsection
\end{frame}
%%%%%%%%%%%%%%%


